%
% TU/e Style Master Thesis template for LaTeX
%
% Public version 1.0
% 2010 - 2013 Thijs Nugteren and Joos Buijs
% Modified and enhanced by Colin Lambrechts, Willem Sonke en Thom Castermans (2015)
%
% THIS IS THE MAIN FILE (i.e. compile this file, compiling the others directly won't work)
%
\documentclass[a4paper,10pt,twoside]{report}

\sloppy %used this to make sure things respect the margin

%all the other includes etc. are done in the thesis.sty file.
\usepackage{thesis}
\usepackage{amsmath,amssymb,amsthm}
\usepackage{wrapfig}
\usepackage{algorithmicx,algorithm,algpseudocode}
\usepackage{nameref}
\usepackage{attachfile}
\usepackage[hang,sf]{caption}
%\usepackage{subcaption}%removed, because it clashes with subfig package
\usepackage{color,soul}

%my favorite font
\usepackage{mathpazo}
\linespread{1.05}

%\usepackage{Alegreya}

%\renewcommand{\familydefault}{\sfdefault}
\usepackage[T1]{fontenc}
%\renewcommand{\rmdefault}{zmb} TUE Meta font

%\usepackage{tikz}%already defined in the .sty
\tikzset{point/.style={circle, fill=black, draw=black, inner sep=2pt}}

\makeatletter
\let\scshape\relax % to avoid a warning
\DeclareRobustCommand\scshape{%
  \not@math@alphabet\scshape\relax
  \ifnum\pdf@strcmp{\f@family}{\familydefault}=\z@
    \fontfamily{jkpkosn}%
  \fi
  \fontshape\scdefault\selectfont}
\makeatother

%set graphics path
\graphicspath{{./figures/}}

%----------------------- Macros and Definitions --------------------------
%line without a number in an algorithm
\def\NoNumber#1{{\def\alglinenumber##1{}\State #1}\addtocounter{ALG@line}{-1}}

%rename Require and Ensure to Input and Output
\floatname{algorithm}{Algorithm}
\renewcommand{\algorithmicrequire}{\textbf{Input:}}
\renewcommand{\algorithmicensure}{\textbf{Output:}}
\newcommand{\from}{\textbf{ from }}
\newcommand{\upto}{\textbf{ up to }}
\newcommand{\downto}{\textbf{ down to }}
\newcommand{\newComment}[1]{\Statex \Comment{#1}}%\strut \\ \Comment{#1}}

%added function labelname
\newfloat{function}{htb}{lop}
\floatname{function}{Function}
\makeatletter
\newcommand{\newalgname}[1]{%
  \renewcommand{\ALG@name}{#1}%
}
%unknown why these two are here
%\newalgname{Algorithm}
%\renewcommand{\listalgorithmname}{List of \ALG@name s}
\makeatother


\newtheorem{definition}{Definition}

\newcommand{\R}{{\mathbb R}}
\newcommand{\dir}{\vec{d}}
\newcommand{\pol}{{\cal P}}
\newcommand{\robot}{{\cal R}}
\newcommand{\triang}{{\cal T}}
\newcommand{\dist}{\mathrm{dist}}
\newcommand{\eps}{\varepsilon}
\newcommand{\Mid}{:}
\newcommand{\ch}{\mathcal{CH}}
\def\qed{\hfill$\Box$}

\newtheorem{theorem}{Theorem}
\newtheorem{lemma}[theorem]{Lemma}

%define TODO with highlighting
\newcommand{\TODO}[1]{\hl{TODO: #1}}


%==========End of self defined macros and definitions =====================
%for fixed table widths
\usepackage{array}
\newcolumntype{L}[1]{>{\raggedright\let\newline\\\arraybackslash\hspace{0pt}}m{#1}}
\newcolumntype{C}[1]{>{\centering\let\newline\\\arraybackslash\hspace{0pt}}m{#1}}
\newcolumntype{R}[1]{>{\raggedleft\let\newline\\\arraybackslash\hspace{0pt}}m{#1}}

%degree symbol
\newcommand{\degree}{^{\circ}}

%-----------------Notions and definitions----------------------------


%
% These commands need to be defined in order to produce a correct and personalized document
%
\newcommand{\shortdoctitle}{short title}
\newcommand{\doctitle}{Title}
\newcommand{\docsubtitle}{Master Thesis}

\newcommand{\me}{Author}
\newcommand{\keywords}{templete, TU, University of Technology}
\newcommand{\version}{Version 0.0}
\newcommand{\monthYear}{\today}

%Be sure to use all the titles for your committee members!!! (their names show up on the very first page!)
\newcommand{\firstCommitteeMember}{\supervisor}
\newcommand{\secondCommitteeMember}{Second committee member}
\newcommand{\thirdCommitteeMember}{Third committee member}
\newcommand{\supervisor}{Supervisor}

\author{\me}

%
% PDF settings
%
\hypersetup
{
    pdfauthor={\me},
    pdftitle={\shortdoctitle},
    pdfsubject={\doctitle},
    pdfkeywords={\keywords},
    linktoc=page
}

\begin{document}

%use this include for PDF and distribution versions
\pagenumbering{roman}
%\include{titlepage}
\newcommand*{\fontscale}{0.85}

\usetikzlibrary{calc}
\tikzset{every picture/.style = {tuedarkpink, line width = 0.25pt}}

% textpos setup, work in mm
\setlength{\TPHorizModule}{1mm}
\setlength{\TPVertModule}{1mm}

% slightly larger linespread
\linespread{1.05}

% titlepage
\begin{titlepage}
%\thispagestyle{empty}

%here come the title
\begin{textblock}{210}(0,103)
  \centering
  \color{tuedarkblue}
  \textbf{\Huge \hspace*{2.3cm} Title
  \\ \vspace*{2.0cm} \Large \hspace*{2.3cm} Master Thesis}
\end{textblock}

%and here the author
\begin{textblock}{210}(0,230)
  \centering
  %\color{tuedarkpink}
  \hspace*{2.3cm} \Large Author

\normalsize
\strut \\ 
\hspace*{2.3cm} Committee:\\
\hspace*{2.3cm} \begin{tabular}{rl}
    \firstCommitteeMember\\
    \secondCommitteeMember\\
    \thirdCommitteeMember\\
\end{tabular}

\strut \\ 
\hspace*{2.3cm} Supervisor:\\
\hspace*{2.3cm} \begin{tabular}{rl}
    \supervisor\\
\end{tabular}

\strut \\
\vfill \hspace*{2.3cm}
%\docdate \\
\large \version \ (\monthYear)\\

%here comes the date of your defense or the publication date
 \strut \\
\vfill \hspace*{2.3cm}
%\docdate \\
\large Eindhoven, \formatdate{9}{9}{2015}\\
\end{textblock}
\color{black}

\begin{tikzpicture}[remember picture, overlay, every text node part/.style = {align = left}]
    \node[inner sep = 0pt, anchor = south east] (naam) at (current page.south east) {%include graphics thing
    \includegraphics[height=29.7cm]{coveredit3}};
  % draw axes
  \draw ($(current page.north east) - (11.7cm, 0)$) -- ++(-106:300);
  \draw ($(current page.north west) - (0, 8.6cm)$) -- ($(current page.north east) - (0, 8.6cm)$);
  % for debugging
  %\draw ($(current page.north west) - (0.1cm, 8.6cm)$) rectangle ($(current page.north east) - (4cm, 14cm)$);
  % fill side
  \begin{scope}
    \clip ($(current page.north west) - (0, 8.6cm)$) rectangle (current page.south east);
    \fill (current page.north west) -- ($(current page.north east) - (11.7cm, 0)$) -- ++(-106:300) -- (current page.south west) -- cycle;
  \end{scope}
  % add logo
  \node[inner sep = 0pt, anchor = north] at ($(current page.north) + (-0.174cm, -1.1cm)$) {\includegraphics[width = 7cm]{tuelogo}};
  % add department
  \node[inner sep = 0pt, anchor = south west, text width = 3.8cm] at ($(current page.north west) + (1.7cm, -8.3cm)$)
    {%\tuedepartmentfont
    \color{tuedarkblue}
    \bfseries\baselineskip=16pt \large / Department of Mathematics and Computer Science \\ / Algorithms \par};
  % add slogan
  \node[inner sep = 0pt, anchor = south west, text width = 3.7cm] at ($(current page.south west) + (1.4cm, 1cm)$)
    {%\tuesloganfont
    \color{tuedarkblue} \footnotesize
    \bfseries Where innovation starts};
\end{tikzpicture}

\end{titlepage}
%~\cleardoublepage

\normalsize

\clearpage

\footnotetext{
Image on the front cover is from ...
}

\clearemptydoublepage

%Sometimes line numbers are nice, uncomment the next line to enable:
%\linenumbers

%It could be handy to have a list of todos and brainstorms in your thesis
%\chapter*{General todos}
%\input{chapters/general_todos}
%TODO maybe change font for dyslectic people?

%\chapter*{*Brainstorm results*}\todo{remove this chapter}
%\input{chapters/brainstorm_results}

\chapter*{Abstract}\label{chapter:abstract}
\input{chapters/abstract}

%used to print the values of the margins
%\printinunitsof{mm}{\pagevalues}

%\clearemptydoublepage
\clearpage %place this here instaed of above line to make sure TOC is on two pages next to each other (starts left and end right)
%note all other chapters start on the right page (just like abstract)

%An executive summary if you want:
%\chapter*{Executive summary}\label{chapter:executive_summary}
%\input{chapters/executive_summary}

%\clearemptydoublepage

%\chapter*{Preface}\label{chapter:preface} %or is it perface
%\input{chapters/preface}

%\clearemptydoublepage

\tableofcontents

\clearemptydoublepage

%TODO: do we need one => no
%\listoffigures

%\clearemptydoublepage

%\listoftables

%\clearemptydoublepage

%\lstlistoflistings

%\clearemptydoublepage

\chapter{Introduction}\label{chapter:Introduction}
\setcounter{page}{0}
\pagenumbering{arabic}
%from here on, start the 'real' page numbering, from 1, with normal digits
\input{chapters/introduction}

\clearemptydoublepage

%probably not needed (only needed when you need an introduction to the study area)
%\chapter{Preliminaries}\label{chapter:preliminaries}
%\input{chapters/preliminaries}
%\clearemptydoublepage

\chapter{Literature and Related Work}\label{chapter:Literature}
\input{chapters/chap_Literature}
\clearemptydoublepage

\chapter{The 1-Dimensional Problem}\label{chapter:Alg1D}
\input{chapters/chap_1D}
\clearemptydoublepage

\chapter{The 2-Dimensional Problem}\label{chapter:Alg2D}
\input{chapters/chap_2D}
\clearemptydoublepage

\chapter{Experiments}\label{chapter:Experiments}
\input{chapters/chap_Experiments}
\clearemptydoublepage

\chapter{The Results and Discussion}\label{chapter:Results}
\input{chapters/chap_Results}
\clearemptydoublepage

\chapter{Conclusions}\label{chapter:Conclusions}
\input{chapters/conclusions}
\clearemptydoublepage

\chapter{Future Work}\label{chapter:FutureWork}
\input{chapters/future_work}

\clearemptydoublepage


%\afterpage{\blankpage} %place the bibliography on the left side, because then all references are visible at the same time
%Choose a good bibliography style, plain would do often, but these might be nice too
%\bibliographystyle{these}
\bibliographystyle{is-abbrv}
\bibliography{references}

%\clearemptydoublepage
\clearpage
%\afterpage{\blankpage}

\phantomsection
\addcontentsline{toc}{chapter}{Appendix}
\appendix

\chapter{Appendix 1 Decisions} \label{chapter:App:1}
\input{appendices/appendix_1}

\end{document}
